\begin{proof}
Let
\[
P:=\operatorname{conv}(S)
\qquad\text{and}\qquad
Q:=\operatorname{conv}(T).
\]
Write the chosen chain in boundary order as
\[
C=(v_1,v_2,\dots,v_m),
\]
and assume that \(m\ge 4\).

By Subproblem~2, \(C\) has an \emph{interior} sign change. Hence there exists an index
\[
i\in\{2,\dots,m-2\}
\]
such that
\[
\varepsilon(t_{v_i})\neq \varepsilon(t_{v_{i+1}}).
\]
In particular,
\[
t_{v_i}\neq t_{v_{i+1}}.
\]

We shall use the standard exposed-face description of Minkowski sums: for every outer normal \(u\),
\[
F_{P+Q}(u)=F_P(u)+F_Q(u),
\]
and therefore
\[
N_{P+Q}(F+G)=N_P(F)\cap N_Q(G).
\]
Equivalently, the normal fan \(\mathcal N(P+Q)\) is the common refinement of \(\mathcal N(P)\) and \(\mathcal N(Q)\), and every vertex of \(P+Q\) decomposes uniquely as a sum of a vertex of \(P\) and a vertex of \(Q\). ([cccg.ca](https://www.cccg.ca/proceedings/2005/44.pdf))

Choose generic directions
\[
u_i\in \operatorname{relint}N_P(v_i),
\qquad
u_{i+1}\in \operatorname{relint}N_P(v_{i+1}).
\]
Then
\[
F_P(u_i)=\{v_i\},
\qquad
F_P(u_{i+1})=\{v_{i+1}\}.
\]
Hence
\[
F_{P+Q}(u_i)=\{v_i+q_i\},
\qquad
F_{P+Q}(u_{i+1})=\{v_{i+1}+q_{i+1}\}
\]
for some vertices \(q_i,q_{i+1}\) of \(Q\). By Subproblem~1, item~\(6\), these two vertices of \(P+Q\) belong to \(\operatorname{supp}(F)\), so by the equality-case assumptions they are lonely points. Since
\[
v_i+q_i\in v_i+T
\qquad\text{and}\qquad
v_{i+1}+q_{i+1}\in v_{i+1}+T,
\]
the uniqueness of the lonely point attached to a given \(s\in S\) yields
\[
v_i+q_i=\ell_{v_i}=v_i+t_{v_i},
\qquad
v_{i+1}+q_{i+1}=\ell_{v_{i+1}}=v_{i+1}+t_{v_{i+1}}.
\]
Therefore
\[
q_i=t_{v_i},
\qquad
q_{i+1}=t_{v_{i+1}}.
\]

Now let \(I\) be the open arc of outer normal directions swept when one moves from \(u_i\) to \(u_{i+1}\) along the chain \(C\). By the monotone normal-sweep statement from the hint, because
\[
t_{v_i}\neq t_{v_{i+1}},
\]
the exposed vertex of \(Q\) changes at some direction
\[
\eta\in I.
\]
Thus the restriction of \(\mathcal N(Q)\) to \(I\) has an interior wall. Since \(\mathcal N(P+Q)\) is the common refinement of \(\mathcal N(P)\) and \(\mathcal N(Q)\), the restriction of \(\mathcal N(P+Q)\) to \(I\) is a \emph{strict} refinement of the restriction of \(\mathcal N(P)\) to \(I\). ([cccg.ca](https://www.cccg.ca/proceedings/2005/44.pdf))

For a polygon, maximal cones of the normal fan correspond, in cyclic order, to boundary vertices. Hence the boundary arc of \(P+Q\) corresponding to \(I\) contains at least one vertex
\[
w
\]
strictly between the endpoint vertices
\[
\ell_{v_i}=v_i+t_{v_i}
\qquad\text{and}\qquad
\ell_{v_{i+1}}=v_{i+1}+t_{v_{i+1}}.
\]

Again by Subproblem~1, item~\(6\),
\[
w\in \operatorname{supp}(F).
\]
We claim that \(w\) is \emph{not} one of the \(n\) lonely points \(\{\ell_s:s\in S\}\).

Indeed, by the unique-decomposition statement above, there are unique vertices \(p\) of \(P\) and \(q\) of \(Q\) such that
\[
w=p+q.
\]
Because the normal cone of \(w\) lies on the boundary arc between the two consecutive chain vertices \(v_i\) and \(v_{i+1}\), the only possible \(P\)-vertices are
\[
p\in\{v_i,v_{i+1}\}.
\]
Since \(w\) is strictly between the endpoints of that arc, it is not equal to
\[
\ell_{v_i}=v_i+t_{v_i}
\quad\text{or}\quad
\ell_{v_{i+1}}=v_{i+1}+t_{v_{i+1}},
\]
so necessarily
\[
q\neq t_p.
\]

Suppose, for contradiction, that
\[
w=\ell_s=s+t_s
\]
for some \(s\in S\). Because every surviving point is lonely, there is exactly one translate \(S+t\) containing \(w\). But
\[
w=p+q\in S+q,
\]
so that unique translate must be \(S+q\). Hence
\[
t_s=q.
\]
Then
\[
s=w-q=p,
\]
and therefore
\[
w=\ell_p=p+t_p,
\]
contradicting \(q=t_s\neq t_p\). This proves that
\[
w\notin \{\ell_s:s\in S\}.
\]

Thus \(w\) is an additional surviving point beyond the \(n\) lonely points, so
\[
|\{x:F(x)\neq 0\}|\ge n+1,
\]
contrary to the equality-case assumption. Therefore no boundary chain can have \(m\ge 4\).

Applying this to both chains gives
\[
|U|\le 3
\qquad\text{and}\qquad
|L|\le 3.
\]
Since \(U\) and \(L\) share exactly their two endpoints,
\[
n=|U|+|L|-2\le 3+3-2=4.
\]
\end{proof}